%%%%%%%%%%%%%%%%%%%%%%%%%%%%%%%%%%%%%%%%%%%%%%%%%%%%%%%%%%%%
% Regulatory Framework
%%%%%%%%%%%%%%%%%%%%%%%%%%%%%%%%%%%%%%%%%%%%%%%%%%%%%%%%%%%%

\section{Regulatory Framework}

India's securitisation market is governed by multiple regulatory bodies to ensure financial stability, investor protection, and transparency. The Reserve Bank of India (RBI) regulates banks and Non-Banking Financial Companies (NBFCs), while the Securities and Exchange Board of India (SEBI) regulates the issuance and trading of securitised securities in the capital market. Credit Rating Agencies (CRISIL, Moody's, S\&P) evaluate the credit quality of securitised instruments.

\vspace{0.4cm}

\subsection*{Key Regulatory Bodies}

\renewcommand{\arraystretch}{1.5}

\begin{tabularx}{\textwidth}{|p{3.5cm}|X|}
	\hline
	\rowcolor{TableHeader}
	\textbf{Organization} & \textbf{Role in Securitisation} \\
	\hline
	
	\textbf{Reserve Bank of India (RBI)}
	&
	Regulates banks and NBFCs involved in securitisation. Issues the Master Directions (2021), defines Minimum Holding Period (MHP), Minimum Retention Requirement (MRR), and risk transfer guidelines.
	\\
	\hline
	
	\textbf{SEBI}
	&
	Regulates the issuance, disclosure, listing, and trading of securitised debt instruments to protect investors.
	\\
	\hline
	
	\textbf{Credit Rating Agencies}
	&
	CRISIL, Moody's, and S\&P assess the credit quality of securitised assets by estimating default risk and expected losses.
	\\
	\hline
	
	\textbf{Banks / NBFCs}
	&
	Originate loans and transfer the loan pools to a Special Purpose Vehicle (SPV) for securitisation.
	\\
	\hline
	
	\textbf{Special Purpose Vehicle (SPV)}
	&
	Purchases loan pools from originators and issues Asset-Backed Securities (ABS) to investors.
	\\
	\hline
	
	\textbf{Investors}
	&
	Purchase securitised securities and receive periodic cash flows generated from borrower repayments.
	\\
	\hline
	
\end{tabularx}

\vspace{0.7cm}

\subsection*{Regulatory Framework Map}

\begin{center}
	
	\fbox{
		\begin{minipage}{0.95\textwidth}
			
			\centering
			
			\textbf{Borrowers}
			
			$\downarrow$
			
			Monthly Loan Repayments
			
			$\downarrow$
			
			\textbf{Banks / NBFCs (Originators)}
			
			$\downarrow$
			
			Transfer Loan Pool
			
			$\downarrow$
			
			\textbf{Special Purpose Vehicle (SPV)}
			
			$\downarrow$
			
			Issues Asset-Backed Securities (ABS)
			
			$\downarrow$
			
			\textbf{Institutional Investors}
			
			\vspace{0.4cm}
			
			\hrule
			
			\vspace{0.2cm}
			
			\textbf{Regulators}
			
			RBI → Banks \& NBFCs
			
			SEBI → Securities Market
			
			Credit Rating Agencies → Credit Assessment
			
		\end{minipage}
	}
	
\end{center}

\vspace{0.5cm}

\subsection*{Key RBI Guidelines (September 2021)}

\begin{itemize}
	
	\item Minimum Holding Period (MHP) before loans can be securitised.
	
	\item Minimum Retention Requirement (MRR) to ensure originators retain part of the credit risk.
	
	\item True Sale Criteria for transferring assets to an SPV.
	
	\item Standardised disclosure requirements for investors.
	
	\item Proper risk transfer from banks to investors.
	
\end{itemize}